\documentclass[11pt,a4paper]{article}

\usepackage[T1]{fontenc}
\usepackage[utf8]{inputenc}
\usepackage[spanish]{babel}
\usepackage{geometry}
\usepackage{microtype}
\usepackage{setspace}
\usepackage{booktabs}
\usepackage{tabularx}
\usepackage{enumitem}
\usepackage{hyperref}
\usepackage{xcolor}
\usepackage{titlesec}
\usepackage{fancyhdr}
\usepackage{tikz}
\usepackage{array}
\usepackage{parskip}

\usetikzlibrary{arrows.meta,positioning}

\geometry{margin=2.3cm, top=2.6cm, bottom=2.6cm}
\setstretch{1.15}
\setlist[itemize]{leftmargin=1.4em,itemsep=0.28em,topsep=0.4em}
\setlist[enumerate]{leftmargin=1.4em,itemsep=0.28em,topsep=0.4em}

\definecolor{upnavy}{HTML}{1F3A5F}
\definecolor{upgray}{HTML}{5A6570}
\definecolor{upsoft}{HTML}{E8EEF4}

\hypersetup{colorlinks=true,urlcolor=upnavy,linkcolor=upnavy}
\titleformat{\section}{\large\bfseries\color{upnavy}}{\thesection.}{0.6em}{}
\titleformat{\subsection}{\normalsize\bfseries\color{upnavy}}{\thesubsection.}{0.5em}{}
\titleformat{\subsubsection}{\normalsize\itshape\color{upnavy}}{\thesubsubsection.}{0.5em}{}

\pagestyle{fancy}
\fancyhf{}
\fancyhead[L]{\small\color{upgray}Reorganización de la carpeta Minería}
\fancyhead[R]{\small\color{upgray}Universidad del Pacífico}
\fancyfoot[C]{\thepage}
\renewcommand{\headrulewidth}{0.3pt}

\newcolumntype{L}[1]{>{\raggedright\arraybackslash}p{#1}}
\newcolumntype{Y}{>{\raggedright\arraybackslash}X}

\begin{document}

\begin{titlepage}
\centering
\vspace*{1.6cm}
{\color{upnavy}\rule{\textwidth}{1.4pt}}\\[0.7em]
{\LARGE\bfseries\color{upnavy} Reorganización del archivo de investigación}\\[0.45em]
{\Large Carpeta Minería --- disco Johao\_asistente}\\[0.7em]
{\color{upnavy}\rule{\textwidth}{1.4pt}}\\[1.4em]
{\large Informe para reunión de seguimiento}\\[2.2em]

\begin{flushleft}
\renewcommand{\arraystretch}{1.35}
\begin{tabular}{@{}l@{\hspace{1.2em}}l@{}}
\textbf{Proyecto} & Minería, conflictos sociales e indicadores territoriales \\
\textbf{Institución} & Universidad del Pacífico \\
\textbf{Disco} & \texttt{Johao\_asistente} \\
\textbf{Carpeta} & Minería \\
\textbf{Preparado por} & Johao Rafael Mendoza Fabián \\
\textbf{Fecha} & Agosto 2026 \\
\end{tabular}
\end{flushleft}

\vfill
\begin{flushleft}
\color{upgray}\small
Este documento explica, con más detalle que la nota corta, por qué se reordenó la carpeta, con qué criterios, cómo quedó el árbol y cómo debería usarse de ahora en adelante. El contenido de investigación no se alteró: se movió, se etiquetó y se documentó.
\end{flushleft}
\end{titlepage}

\tableofcontents
\newpage

\section{Para qué sirve este informe}

La carpeta Minería acumuló, a lo largo del proyecto, datos de varias fuentes, notebooks, mapas, papers y entregas. Con el tiempo eso se volvió difícil de leer: no había un lugar obvio por donde empezar, los nombres no decían el rol de cada cosa y había copias del mismo material en más de un sitio.

El trabajo que se hizo fue de archivo, no de análisis. No se reestimó ningún indicador ni se reescribió código de investigación. Lo que cambió es el mapa: dónde está cada tipo de archivo, qué se considera vigente y qué se guarda por si hace falta volver atrás.

Este informe está pensado para la reunión: se puede recorrer en orden, o saltar a la sección que interese (datos, código, archivo, reglas del equipo). Al final hay un árbol completo y una guía de diez minutos para alguien que entre por primera vez.

\section{De qué va el proyecto, en breve}

Minería es el archivo de un trabajo de la Universidad del Pacífico sobre minas en Perú. La pregunta de fondo es cómo se relaciona el tipo de operación minera ---tamaño, método, formalidad--- con los conflictos sociales (económicos, ambientales, políticos) y con lo que ocurre en los pueblos del entorno: población, salud, educación, infraestructura. El radio de trabajo que aparece en la bitácora es de unos 20~km alrededor de la mina.

Las operaciones que ya tienen material propio en la carpeta son, entre otras, Toromocho (Junín, Chinalco), Quellaveco (Moquegua, Anglo American), Yanacocha y Antamina. En las notas internas se habla de ampliar a unas diez minas, con cobertura regional, y de un universo más grande de minas formales en producción.

Para esa pregunta el archivo junta tres familias de evidencia:

\begin{itemize}
  \item \textbf{Territorio e imágenes.} Landsat y Sentinel, índices de vegetación y humedad (NDVI, NDWI, NDMI), shapefiles del IGN y capas de área en operación.
  \item \textbf{Conflictos.} Reportes mensuales de la Defensoría del Pueblo y registros de OCMAL, pasados luego a JSON, Excel y bases de trabajo.
  \item \textbf{Socioeconomía.} PBI subnacional, demografía y otros indicadores de contexto.
\end{itemize}

Eso es lo que había que poder encontrar sin conocer la historia de quién subió cada archivo.

\section{Cómo estaba la carpeta antes}

En la raíz de Minería convivían, sin un orden de lectura, la bitácora, un documento de puntos pendientes, JSON sueltos, Excel, Word, un Colab sin extensión clara, notebooks, papers y una segunda entrega. Las carpetas de primer nivel tenían nombres de trabajo ---Defensoría del Pueblo, Imágenes satelitales, Informe, Mapa de minas, Pasado código--- que tenían sentido para quien las creó, pero no para alguien que llega después.

Tres problemas se repetían:

\begin{enumerate}
  \item \textbf{Mezcla de roles.} En el mismo nivel estaban insumos, código, resultados, lecturas y entregas. No se distinguía qué era fuente, qué era producto y qué era copia de seguridad.
  \item \textbf{Duplicados.} Los PDF de la Defensoría aparecían en más de una carpeta (del orden de 250 reportes en la pila principal, con otra pila equivalente). Había series NDVI/NDWI repetidas y más de una copia del notebook de indicadores.
  \item \textbf{Nombres poco estables.} Variantes como \texttt{Información\_}, \texttt{Segmentation}, \texttt{Mapa\_minas} o carpetas de ``pasado código'' obligaban a abrir varias veces para saber cuál era la versión que se usa.
\end{enumerate}

Un compañero nuevo no podía responder, en dos minutos, tres preguntas básicas: ¿por dónde empiezo?, ¿cuál es el archivo vigente?, ¿dónde dejo algo nuevo sin ensuciar el resto?

\section{Criterios con los que se ordenó}

La reorganización se hizo solo dentro de Minería, en el disco compartido \texttt{Johao\_asistente}. No se tocó la otra copia de Minería que aparece en ``compartido conmigo'' ni los ZIP de imágenes que están en la raíz del disco. Nada se borró de forma definitiva: se movió. Lo duplicado o viejo fue a archivo, no a la papelera como política de trabajo.

Los criterios fueron estos:

\subsection{Un orden de lectura, no un orden alfabético}

Drive ordena por nombre. Si las carpetas se llaman ``Análisis'', ``Código'', ``Datos'', el orden alfabético no coincide con cómo se entra a un proyecto. Por eso van numeradas: \texttt{00} se lee primero, \texttt{99} se deja para el final. Quien abre la carpeta ve de arriba hacia abajo el flujo natural: entender, datos, código, resultados, lecturas, entregas, archivo.

\subsection{Separar lo que no se edita de lo que sí}

Los PDF de la Defensoría, el JSON crudo de OCMAL, los shapefiles y las series oficiales de PBI o demografía son insumos. Si alguien los pisa, se pierde la trazabilidad. Las tablas y bases que el equipo ya armó ---conflictos, minas, indicadores satelitales--- son productos. Van en cajones distintos: \texttt{crudos} y \texttt{procesados}.

\subsection{Agrupar el código por tema, no por la carpeta en la que nació}

Había notebooks repartidos en ``Informe'', ``Imágenes'' y ``Pasado código''. Ahora el código vigente está por pregunta: conflictos de Defensoría, OCMAL, imágenes satelitales, mapas, segmentación. Si alguien busca ``el notebook de NDVI'', va a una sola carpeta.

\subsection{No borrar; archivar}

Un archivo de investigación se equivoca hacia conservar. Lo que estaba repetido o pertenecía a un pipeline viejo (por ejemplo extracción 2004--2014) pasó a \texttt{99\_Archivo}, con la palabra \texttt{DUPLICADO} en el nombre cuando hacía falta. El vigente se busca en \texttt{01} o \texttt{02}. Si mañana hace falta comparar, el material sigue ahí.

\subsection{Dejar pistas escritas}

Cada carpeta grande tiene un README corto, en español, como nota de equipo. En datos hay un diccionario. En inicio hay cómo usar la carpeta y convenciones de nombres. La idea es que no haga falta haber estado en las reuniones previas para orientarse.

\section{Mapa general: qué hay ahora en la raíz}

En la raíz de Minería quedan únicamente el índice \texttt{LEEME.md} y siete carpetas. Nada suelto.

\begin{center}
\renewcommand{\arraystretch}{1.25}
\begin{tabularx}{\textwidth}{@{}c l Y@{}}
\toprule
\textbf{N.º} & \textbf{Carpeta} & \textbf{Función} \\
\midrule
00 & \texttt{00\_Inicio} & Qué es el proyecto, bitácora, pendientes y reglas de uso. \\
01 & \texttt{01\_Datos} & Fuentes crudas y bases ya armadas. \\
02 & \texttt{02\_Codigo} & Notebooks vigentes, agrupados por tema. \\
03 & \texttt{03\_Analisis} & Mapas HTML y fichas de minas. \\
04 & \texttt{04\_Literatura} & Papers y material del seminario. \\
05 & \texttt{05\_Entregas} & Snapshots de lo ya enviado (p.~ej.\ segunda entrega). \\
99 & \texttt{99\_Archivo} & Copias, código viejo y duplicados. \\
\bottomrule
\end{tabularx}
\end{center}

\vspace{0.8em}
El flujo que se espera al trabajar es el de la figura: se lee el inicio, se toman datos, se corre código, se guardan resultados; las lecturas y las entregas viven aparte; el archivo no se usa en el día a día.

\begin{center}
\begin{tikzpicture}[
  font=\small,
  box/.style={draw=upnavy, rounded corners=2pt, align=center, inner sep=6pt, fill=upsoft, text=upnavy, minimum height=1.05cm, text width=2.15cm},
  arr/.style={-{Latex}, thick, color=upnavy}
]
\node[box] (a) {00\\Inicio};
\node[box, right=0.45cm of a] (b) {01\\Datos};
\node[box, right=0.45cm of b] (c) {02\\Código};
\node[box, right=0.45cm of c] (d) {03\\Análisis};
\node[box, below=0.85cm of b] (e) {04\\Literatura};
\node[box, below=0.85cm of c] (f) {05\\Entregas};
\node[box, below=0.85cm of d] (g) {99\\Archivo};
\draw[arr] (a) -- (b);
\draw[arr] (b) -- (c);
\draw[arr] (c) -- (d);
\draw[arr] (a) -- (e);
\draw[arr] (c) -- (f);
\draw[arr] (b) -- (g);
\end{tikzpicture}
\end{center}

\section{Detalle por carpeta}

\subsection{\texttt{00\_Inicio} --- por dónde se entra}

Aquí está lo que antes estaba suelto en la raíz y lo que se escribió para que el archivo se explique solo:

\begin{itemize}
  \item un README de la carpeta;
  \item \texttt{COMO\_USAR}: si uno viene por datos, por código o por el estado del avance;
  \item \texttt{CONVENCIONES}: cómo nombrar y dónde guardar cosas nuevas;
  \item \texttt{Bitacora\_de\_avance.docx};
  \item \texttt{Puntos\_a\_considerar.docx}.
\end{itemize}

La bitácora y los pendientes no se reescribieron: solo cambiaron de sitio, para que el estado del proyecto viva junto a las reglas de la carpeta y no entre JSON y Excel.

\subsection{\texttt{01\_Datos} --- crudos y procesados}

Es la carpeta más importante para no perder trazabilidad. Tiene dos subcarpetas grandes y un \texttt{DICCIONARIO\_DATOS.md} que dice qué es cada JSON o Excel.

\subsubsection{Crudos}

Lo que vino de afuera y no conviene editar:

\begin{center}
\renewcommand{\arraystretch}{1.2}
\begin{tabularx}{\textwidth}{@{}l Y@{}}
\toprule
\textbf{Subcarpeta} & \textbf{Qué hay} \\
\midrule
\texttt{defensoria\_pdfs} & Reportes mensuales de conflictos de la Defensoría del Pueblo (la pila vigente; hay una carpeta \texttt{PDFS} anidada con esos archivos). \\
\texttt{ocmal} & JSON crudo de OCMAL. \\
\texttt{espacial} & Shapefiles del IGN (\texttt{peru\_shapes}) y geojson de área en operación. \\
\texttt{socioeconomico} & DI20, PBI subnacional, evolución demográfica. \\
\bottomrule
\end{tabularx}
\end{center}

\subsubsection{Procesados}

Lo que el equipo ya dejó listo para analizar:

\begin{center}
\renewcommand{\arraystretch}{1.2}
\begin{tabularx}{\textwidth}{@{}l Y@{}}
\toprule
\textbf{Subcarpeta} & \textbf{Qué hay} \\
\midrule
\texttt{conflictos} & \texttt{conflictos\_data.json}, base de Defensoría, pueblos, Excel de OCMAL. \\
\texttt{minas} & \texttt{minas\_df}, unidades minerales, listados de trabajo. \\
\texttt{indicadores\_satelitales} & \texttt{INDICADORES\_MINAS} y series NDVI, NDWI, NDMI. \\
\bottomrule
\end{tabularx}
\end{center}

Regla simple: si alguien extrae una tabla nueva, va a procesados. Los crudos se consultan, no se pisan. Si hay duda de qué archivo usar, el diccionario de esa misma carpeta es la referencia, no el recuerdo de quién lo subió.

\subsection{\texttt{02\_Codigo} --- notebooks por tema}

El código vigente quedó agrupado así:

\begin{itemize}
  \item \texttt{conflictos\_defensoria}: extracción y armado a partir de los reportes de la Defensoría.
  \item \texttt{ocmal}: scraping y limpieza de OCMAL.
  \item \texttt{imagenes\_satelitales}: índices y series (NDVI, NDWI, NDMI).
  \item \texttt{mapas}: construcción de los HTML de minas.
  \item \texttt{segmentacion}: notebooks de segmentación de imágenes.
\end{itemize}

Cada una de esas carpetas concentra el notebook que se usa ahora. El código viejo ---por ejemplo el pipeline de extracción 2004--2014--- no está aquí: está en archivo, para no confundirlo con la versión actual.

\subsection{\texttt{03\_Analisis} --- lo que se muestra}

Aquí van productos que se abren o se leen como resultado, no como insumo:

\begin{itemize}
  \item mapas en HTML (para el navegador);
  \item fichas de minas (Toromocho, Quellaveco, Antamina, y las que se vayan sumando).
\end{itemize}

No están mezclados con papers ni con los PDF de la Defensoría. Si mañana hay un mapa nuevo, este es el lugar, no la raíz ni \texttt{01\_Datos}.

\subsection{\texttt{04\_Literatura}}

Papers de teledetección, conflictos y el material del seminario de la UP. Separar lecturas de datos evita que un PDF de un paper se pierda entre doscientos reportes mensuales de la Defensoría.

\subsection{\texttt{05\_Entregas}}

Lo que ya se mandó queda como snapshot. Hoy está, por ejemplo, \texttt{2023-10\_segunda\_entrega} (listado Huancayo e imágenes de esa fecha). No se mezcla con lo que se sigue editando: si hay que reconstruir qué se entregó, se abre esta carpeta, no se adivina entre versiones sueltas.

\subsection{\texttt{99\_Archivo}}

Borradores, el pipeline viejo, la segunda pila de PDF de Defensoría, copias de Excel y notebooks, series NDVI/NDWI duplicadas. Cuando el nombre no bastaba, se añadió \texttt{DUPLICADO}. Esta carpeta no es basura: es el lugar donde se guarda lo que no debe competir con el archivo vigente.

\section{Qué se movió, en la práctica}

Antes de tener permiso de administración en el disco compartido, la cuenta de trabajo podía crear carpetas y documentación, pero no reubicar archivos de otras personas. En esa etapa se llegó a usar atajos. Cuando el rol de administrador de contenido quedó activo, los atajos se retiraron y los archivos reales se movieron a \texttt{00}--\texttt{99}. Las carpetas vacías del árbol viejo (\texttt{z\_legado\_*}) se dejaron de usar.

El resultado que hay que mostrar en la reunión es ese estado final: en la raíz, solo el índice y las siete carpetas numeradas; adentro, los archivos de verdad, no un segundo árbol paralelo.

Un detalle menor que quedó del árbol original: los PDF de la Defensoría viven en \texttt{01\_Datos/crudos/defensoria\_pdfs/PDFS} (una carpeta extra). No afecta el uso; es el nombre que ya traían.

\section{Documentación que se dejó en Drive}

Además de mover archivos, se escribieron notas cortas en cada carpeta grande. El tono es de equipo, no de manual.

\begin{center}
\renewcommand{\arraystretch}{1.2}
\begin{tabularx}{\textwidth}{@{}l Y@{}}
\toprule
\textbf{Archivo} & \textbf{Para qué} \\
\midrule
\texttt{LEEME.md} (raíz) & Puerta de entrada: las siete carpetas en orden. \\
\texttt{00\_Inicio/README.md} & De qué va el trabajo y qué hay en inicio. \\
\texttt{COMO\_USAR} & Rutas distintas según si uno busca datos, código o el avance. \\
\texttt{CONVENCIONES} & Nombres y destinos de archivos nuevos. \\
\texttt{01\_Datos/README.md} & Crudos frente a procesados. \\
\texttt{DICCIONARIO\_DATOS.md} & Qué es cada JSON, Excel o shapefile. \\
README de \texttt{02} a \texttt{99} & Qué hay en código, análisis, literatura, entregas y archivo. \\
\bottomrule
\end{tabularx}
\end{center}

La bitácora y los puntos pendientes siguen siendo los documentos Word del equipo. La documentación nueva no los reemplaza: los ubica.

\section{Cómo recorre esto alguien que acaba de entrar}

Una ruta razonable, de unos diez minutos:

\begin{enumerate}
  \item Abrir \texttt{LEEME.md} en la raíz. Ahí está el mapa de una página.
  \item Entrar a \texttt{00\_Inicio}: leer el README, mirar la bitácora y los puntos pendientes.
  \item Si el encargo es de datos, ir a \texttt{01\_Datos}, leer el diccionario y trabajar sobre \texttt{procesados}; consultar \texttt{crudos} solo para volver a la fuente.
  \item Si el encargo es de código, ir a \texttt{02\_Codigo} y abrir la subcarpeta del tema (no \texttt{99\_Archivo}).
  \item Si hay que ver un mapa o una ficha, \texttt{03\_Analisis}.
  \item Si hay que citar papers o material del seminario, \texttt{04\_Literatura}.
  \item Si hay que saber qué se entregó en una fecha, \texttt{05\_Entregas}.
\end{enumerate}

Con eso no hace falta preguntar ``¿en qué carpeta estaba el Excel de OCMAL?''. La respuesta está en el diccionario y en procesados/conflictos.

\section{Qué se pide al equipo para que no se desarme}

El orden se mantiene si se respetan cuatro hábitos:

\begin{enumerate}
  \item \textbf{Entrar por el índice.} \texttt{LEEME.md} y \texttt{00\_Inicio}, no un archivo suelto en la raíz.
  \item \textbf{No dejar huérfanos.} Nada nuevo en la raíz de Minería. Todo va a una de las siete carpetas.
  \item \textbf{No editar crudos.} Las salidas nuevas van a \texttt{procesados} o a \texttt{03\_Analisis}.
  \item \textbf{Las copias, al archivo.} Si se duplica un Excel o un notebook ``por si acaso'', el lugar es \texttt{99\_Archivo}, no al lado de la versión que sí se usa.
\end{enumerate}

Convenciones de nombre que ya quedaron escritas: carpetas numeradas con guion bajo; temas en minúsculas y sin espacios (\texttt{imagenes\_satelitales}, no ``Imágenes Satelitales 2''); entregas con fecha al frente (\texttt{2023-10\_segunda\_entrega}).

\section{Qué no cambió (y conviene decirlo en la reunión)}

\begin{itemize}
  \item El contenido de investigación es el mismo. No se rehicieron indicadores ni se reescribieron notebooks de análisis.
  \item No se tocó la otra carpeta Minería de ``compartido conmigo'' ni los ZIP de la raíz del disco.
  \item No se eliminó material de trabajo: lo duplicado o viejo está en \texttt{99\_Archivo}.
  \item La bitácora y los puntos pendientes siguen siendo los documentos que el equipo ya usaba.
\end{itemize}

En una frase para cerrar: cambió el mapa para encontrar el trabajo; no el trabajo en sí.

\section{Árbol de referencia}

\vspace{0.4em}
{\small\ttfamily
\begin{tabular}{@{}l@{}}
Minería/ \\
\quad LEEME.md \\
\quad 00\_Inicio/ \\
\quad\quad README.md, COMO\_USAR, CONVENCIONES \\
\quad\quad Bitacora\_de\_avance.docx \\
\quad\quad Puntos\_a\_considerar.docx \\
\quad 01\_Datos/ \\
\quad\quad README.md, DICCIONARIO\_DATOS.md \\
\quad\quad crudos/ \\
\quad\quad\quad defensoria\_pdfs/PDFS \\
\quad\quad\quad ocmal/ \\
\quad\quad\quad espacial/ \\
\quad\quad\quad socioeconomico/ \\
\quad\quad procesados/ \\
\quad\quad\quad conflictos/ \\
\quad\quad\quad minas/ \\
\quad\quad\quad indicadores\_satelitales/ \\
\quad 02\_Codigo/ \\
\quad\quad conflictos\_defensoria/ \\
\quad\quad ocmal/ \\
\quad\quad imagenes\_satelitales/ \\
\quad\quad mapas/ \\
\quad\quad segmentacion/ \\
\quad 03\_Analisis/ \\
\quad\quad mapas\_html/ \\
\quad\quad fichas\_minas/ \\
\quad 04\_Literatura/ \\
\quad\quad teledeteccion/, conflictos/, seminario\_up/ \\
\quad 05\_Entregas/ \\
\quad\quad 2023-10\_segunda\_entrega/ \\
\quad 99\_Archivo/ \\
\quad\quad borradores, pipeline viejo, duplicados (PDF, notebooks, xlsx) \\
\end{tabular}
}

\vspace{1.2em}
\noindent\textit{Enlace de la carpeta:} \url{https://drive.google.com/drive/folders/1NfRO7wAfOVF6dasoElMsV8srW-7ThSh9}

\end{document}
